\chapter{Comparison with Related Geometric Approaches}
\label{ch:ch14_related}

The question of this chapter is where the constructions of this manuscript stand among other attempts to read particle physics off an
algebraic or geometric structure. Comparisons that concern a single statement are made where the statement occurs, and this chapter only
points to them. It then treats in more detail the two programmes closest in method: Krasnov's realisation of one generation on the forms of
a doubled split space, which shares the split signature and the $64$-dimensional generation (\cref{sec:rel:krasnov}), and the spectral
programme of Chamseddine and Connes, which shares the strategy of reading the gauge symmetry off an algebra (\cref{sec:rel:connes}). The last
section separates what is inherited, what is a new realisation, what is a new classification or obstruction, and what would be a prediction
(\cref{sec:rel:novelty}). The comparisons rest on the sources as read and recorded by the author; equal
dimensions are never taken as equivalence.

\section{Related work}
\label{sec:rel:work}

\paragraph{Clifford and division algebras.} The classification of real Clifford algebras and their gradings is classical
\cite{LawsonMichelsohn,Lounesto,Porteous}, and the use of Clifford algebras as a language for relativistic physics goes back to Hestenes'
space-time algebra \cite{Hestenes66,Hestenes67}; the one-equation form of Maxwell's equations used on the octants in \cref{ch:ch03_readings} is
Hestenes' form transferred to another module. Several programmes extract Standard-Model structure from division and Clifford algebras
\cite{Furey,Gresnigt}, and the representation-theoretic limits of embedding all of it in a single exceptional group are discussed in
\cite{DistlerGaribaldi}. This manuscript shares the method of reading internal symmetry off centralisers of a Lorentz action in a Clifford
algebra. It differs in the algebra, the split real algebras $\Cl(3n,3n)$, and in its emphasis: it records which choices the algebra removes and
which it leaves, and treats a fit of Standard-Model content as a consistency result, not as evidence (\cref{cor:nopredict}).

\paragraph{Pati--Salam and $SO(10)$.} The content of the chiral block, one Pati--Salam family with a right-handed neutrino
(\cref{thm:PS}), is that of the Pati--Salam model \cite{PatiSalam} and of the spinor $\mathbf{16}$ of $SO(10)$ \cite{FritzschMinkowski}; its
anomaly freedom is inherited \cite{GeorgiGlashow72}. The reading of the $272$ mass kernels through the $SO(10)$ representations $\mathbf{10}$ and $\mathbf{126}$
(\cref{thm:kernelspace}) and the seesaw and Georgi--Jarlskog structures of \cref{ch:ch07_vacuum} are standard model building. What is specific here is that
they occur inside the realisation $V_*=\operatorname{im}P_*$ and are computed in it.

\paragraph{Gravity and unification in split signature.} Embeddings of one generation together with the Lorentz group into a single orthogonal
group, $SO(3,11)$ and $SO(7,7)$, were proposed by Nesti and Percacci \cite{NestiPercacci10} and by Maraner \cite{Maraner04}, and Krasnov's
construction (\cref{sec:rel:krasnov}) realises the $SO(7,7)$ case on forms. The tower $\Cl(3n,3n)$ at $n=3$ has eighteen generators, not
fourteen, and its matter block is selected by a projector rather than by a Weyl condition, so none of these embeddings is reproduced here;
nor does this manuscript derive gravitational dynamics (\cref{sec:gravlimits}).

\paragraph{Forms as spinors.} The realisation of the spinor module of a split space on the exterior algebra of a null half-space is
Chevalley's \cite{Chevalley54}, and the reading of the operator $d-\delta$ on forms as a Dirac operator is the Dirac--K\"ahler construction
\cite{Kahler62,Graf78,BecherJoos82,BennTucker83}. Both are inherited at one constituent (\cref{r3:prop:fock}); at three constituents the manuscript
uses no form realisation.

\paragraph{Electromagnetism.} The comparisons for the field-strength reading are made in \cref{ch:ch03_readings}, each with its class: the
Chevalley--Mukai pairing and the pure-spinor criterion \cite{Chevalley54,Gualtieri03}, Rainich's identity \cite{Rainich1925}, the ModMax family
\cite{BandosLechnerSorokinTownsend20,Kosyakov20}, and the textbook Pauli and electric-dipole couplings \cite{Pauli41}, whose real forms and count at
every clock are given in \cref{r4:thm:couplings}. The duality between the copy algebra and the diagonal Lorentz algebra along the tower is an
instance of Howe duality \cite{Howe89} (\cref{cor:howe}).

\section{Krasnov's construction}
\label{sec:rel:krasnov}

Krasnov \cite{Krasnov18} realises the spinors of $SO(n,n)$ on the inhomogeneous forms of a null $n$-plane, in the setting of doubled
geometry, where the Dirac operator acts by wedge in one half of the coordinates and by contraction in the other. The Lorentz group is
selected by a second splitting, a choice that the paper does not make dynamical, and all fields depend on four of the $2n$ coordinates. At
$SO(3,3)$ the commutant of the Lorentz group is $SO(2)$ and the kinetic term reduces to that of one charged two-component Weyl fermion. At
$SO(7,7)$ the commutant is $SO(4,6)$, whose compact part $SO(4)\times SO(6)$ is the Pati--Salam group, and one generation with $\nu_R$ is a real
even form on $\R^7$, of real dimension $64$; in the kinetic term the $SU(2)_L$ and $SU(2)_R$ sectors enter with opposite signs. Bosons and
symmetry breaking are left open. \Cref{tab:krasnov} compares the two constructions along the axes that matter.

\begin{table}[t]
\centering
\footnotesize
\begin{tabular}{@{}>{\raggedright\arraybackslash}p{2.4cm}>{\raggedright\arraybackslash}p{4.0cm}>{\raggedright\arraybackslash}p{4.6cm}>{\raggedright\arraybackslash}p{3.6cm}@{}}
\toprule
axis & Krasnov \cite{Krasnov18} & this manuscript & relation\\
\midrule
ambient algebra & $\R^{7,7}$ doubled, $\Cl(7,7)$ & $\R^{3,3}$ and the tower $\Cl(3n,3n)$; at $n=3$, $\Cl(9,9)$ on $\R^{512}$ (\cref{sec:tower}) & analogous (split); different ambient (14 vs 18 generators)\\
module & $\Lambda^\bullet\R^7=\R^{128}$; Weyl half $\Lambda^{\rm even}\R^7=\R^{64}$ & $V_*=\operatorname{im}P_*$, real dimension $64$ (\cref{sec:pstar}) & same dimension; no equivalence shown\\
forms or Clifford module & forms on a null half (Chevalley); $D$ by wedge and contraction & Kronecker realisation; forms used at $n=1$ only (\cref{r3:prop:fock}) & the same Chevalley construction at $n=1$\\
complex structure, chirality & fixed even parity; real Weyl spinor & $J$, the product of the 14 internal generators; $\Gamma=J\Omega$ (\cref{thm:selector3}) & analogous; mechanisms differ\\
Lorentz embedding & a second splitting, not dynamical & the interior $\so{1,3}$ after a clock, imposed (\cref{assum:clockred}) & analogous: both are choices\\
lowest rung & $SO(3,3)$: commutant $SO(2)$, one charged Weyl fermion & $n=1$: internal centraliser $\uu1=\langle J\rangle$, $J$ the imaginary unit of the Dirac module (\cref{prop:6plus8plus1,thm:i}) & equivalent as a Lie-algebra statement; field content differs\\
Pati--Salam rung & $SO(4,6)\supset SO(4)\times SO(6)$ acts on the whole $64$ & $P_*$ preserved by the $21$ compact generators only; each of the $24$ non-compact ones maps $V_*$ to its complement (\cref{rem:krasnovpointer}, z3-9) & not the same action\\
one generation & SM generation with $\nu_R$ & one Pati--Salam family with $\nu_R$ (\cref{thm:PS}) & same count; equivalence as modules not constructed\\
projector & Weyl condition & $P_*=\tfrac18(\id+\mathcal X)(\id+\mathcal Z)(\id+\mathcal U)$, \cref{eq:Pstar} & different mechanisms\\
kinetic term & $SU(2)_L$ and $SU(2)_R$ with opposite signs & $K^\mu=P_*\gamma^0\gamma^\mu P_*$, $K^0=\id$ on $V_*$, one sign & different; comparison needs a dictionary\\
mass bilinears & none for the Weyl Lagrangian & $272$-dimensional kernel space (\cref{thm:kernelspace}) & treated here only\\
bosons and actions & proposed $D\to D+C$; not derived & adopted action \eqref{eq:Saction} and potential & neither derives the bosonic action\\
\bottomrule
\end{tabular}
\caption{Krasnov's construction and the chiral block of this manuscript. ``Equivalent'' is used only where a statement is proved in both; equal
dimensions are not equivalence.}
\label{tab:krasnov}
\end{table}

\paragraph{What is shared and what is not.} Two statements are equivalent: the Chevalley realisation of the spinors of a split space on the forms
of a null half-space, which both constructions inherit, and the commutant $\so2$ of a Lorentz subalgebra in $\so{3,3}$. Three are analogous:
Lorentz and Pati--Salam in one split-signature Clifford algebra with a $64$-dimensional one-generation module; the non-dynamical selection of the
Lorentz subgroup; and the imposed reduction to four coordinates. The two $64$-dimensional modules carry different actions: Krasnov's carries
$\so{4,6}$, the block $V_*$ only its compact part.

\paragraph{No representation map has been constructed.} No intertwiner between $V_*$ and $\Lambda^{\rm even}\R^7$ is written down here, and
none is claimed. An explicit map would need a common embedding of $\so{1,3}\oplus\so4\oplus\so6$, an identification of the relevant
$\Cl(9,9)$ block with $\Cl(7,7)$ that drops or absorbs the extra positive four-plane on which the normal charge $Q_A$ lives, and intertwining
identities for the kinetic kernels, including the relative sign of Krasnov's two sectors. That construction is open.

\section{The spectral programme of Chamseddine and Connes}
\label{sec:rel:connes}

Chamseddine and Connes obtain the Standard Model from a spectral triple. An algebra $\mathcal A=C^\infty(M)\otimes\mathcal A_F$ acts on a
Hilbert space with a Dirac operator $D$, a real structure $J_{\rm C}$ and a grading $\gamma$; the finite algebra $\C\oplus\mathbb H\oplus M_3(\C)$ is
singled out among the irreducible finite geometries of KO-dimension $6$ that satisfy the order-one condition, under a symplectic--unitary
hypothesis, which also fixes $16$ fermions per generation \cite{ChamseddineConnes07,ChamseddineConnes08}. Gauge and Higgs fields arise as inner
fluctuations $D\mapsto D+A+J_{\rm C}AJ_{\rm C}^{-1}$, and the spectral action $\Tr f(D/\Lambda)$ yields, in the Euclidean formulation of
\cite{ChamseddineConnes96,ChamseddineConnes}, the bosonic Standard-Model action coupled to gravity, with the gauge couplings related at the
cut-off scale. Its inputs are the number of generations, the finite Dirac operator with its Yukawa and Majorana data, the cut-off function and
scale, and the hypotheses of the classification. This manuscript shares with that programme the strategy of reading the gauge symmetry off
an algebraic structure, and none of its objects. It has no algebra subject to order-zero and order-one conditions, no real structure of
KO-dimension $6$ and no spectral action: its gauge algebras are centralisers in a split Clifford algebra, its matter block is the image of a
chosen projector, and its actions are adopted, not derived from a spectral principle. Its inputs are the clock, the projector, the Lorentz
action, the actions and potential, and the scales. It is therefore not a spectral-triple theory, and no number of the spectral-action
programme transfers to it. The two meet the normalisation $\sin^2\theta_W=\tfrac38$ only through the grand-unified relation of couplings at a
high scale, which is a normalisation and not a prediction without a matching scale, a spectrum and running (PH-09).

\section{What is inherited and what is new}
\label{sec:rel:novelty}

The statements of this manuscript fall into four kinds, and the comparisons above place them (CORE-13).
\begin{enumerate}[label=(\roman*),leftmargin=2.2em]
\item \emph{Known identities, restated in the arena's variables.} Ordinary complex Dirac theory (\cref{thm:realification}); the Chevalley
realisation; the Lagrange and Rainich identities of the field energy; Maxwell's equations and their duality; the Pati--Salam and $SO(10)$
content and its anomaly freedom; Kaluza--Klein reduction, O'Neill's formulas, Lanczos's identity and the Holst action; Howe duality; the
reparametrisation invariance of a worldline.
\item \emph{New realisations.} The tower $\Cl(3n,3n)$ as a setting for these structures; the chiral projector $P_*$ built from three commuting
involutions, with the normal charge on an extra positive four-plane; the field-strength reading with $G$ as the Hodge star and the duality
phase of the curvature as a Barbero--Immirzi angle; the interior dynamics, its closure and its rigid-body realisation.
\item \emph{New classifications and obstructions, within this realisation.} The internal centralisers along the tower; the $272$-dimensional
kernel space and the vacuum analysis on it; the exclusion of the two readings from one another; the obstructions to hypercharge from anomaly
cancellation in the class of three involutions, to colour in the five-projector family, to the normal $U(1)$ as the photon, and to type
selection by invariant scalar dynamics.
\item \emph{Predictions of a specified model.} None is claimed.
\end{enumerate}
The novelty of the items in (ii) and (iii) is relative to the sources read for this chapter. It rests on a targeted search, not on a survey of
the literature, and no blanket claim of originality follows from it.
